%%%%%%%%%%%%%%%%%%%%%%%%%%%%%%%%%%%%%%%%%
%
% (c) 2018 by Jennifer Laaser
%
% This work is licensed under the Creative Commons Attribution-NonCommercial-ShareAlike 4.0 International License. To view a copy of this license, visit http://creativecommons.org/licenses/by-nc-sa/4.0/ or send a letter to Creative Commons, PO Box 1866, Mountain View, CA 94042, USA.
%
% The current source for these materials is accessible on Github: https://github.com/jlaaser/pogil-polymers
%
%%%%%%%%%%%%%%%%%%%%%%%%%%%%%%%%%%%%%%%%%

\renewcommand{\figpath}{content}
\renewcommand{\labelbase}{labelbase}

\begin{activity}[Activity Template]
\label{\labelbase}

\begin{instructornotes}
	Here is some information for the instructor about doing this activity.
	
	Eventually, the \verb!pogil! class may be updated to include macros to standardize formatting, but for now, you will probably find it useful to include information about:
	
	\begin{itemize}
		\item Activity type
		\item Content goals
		\item Process goals
		\item Estimated time required
		\item Instructor preparation required
		\item Prerequisite knowledge
	\end{itemize}
	
\end{instructornotes}

	\textbf{Focus question:} Put a central question for the students to consider through this exercise here.



\begin{model}[ABC]
\label{\labelbase:mdl:MODEL1}

	Here is the first model for students to consider

	Here is a second paragraph of text, followed by an image:

	% to include images, put them in the folder specified by figpath and then use:
	\centerline{\includegraphics[width=0.5\textwidth]{\figpath/test_image.pdf}}

\end{model}


\begin{ctqs}

	\question First question?  I can reference Model \ref{\labelbase:mdl:MODEL1} using a appropriate ref command.
	
		\begin{solution}[1in]{}
		\end{solution}
		
	\question Second question?
	
		\begin{solution}[1in]{}
			Here is a dummy answer in the solution space.
		\end{solution}
		
	\question Third question?
	
		\begin{solution}[1in]{Here is something we show only when we process the student version.}
			
			Here is what we show when we process the instructor version.
			
		\end{solution}

\end{ctqs}



\begin{model}[DEF]
\label{\labelbase:mdl:MODEL2}

	Here is a second model for students to consider.

\end{model}

\begin{ctqs}

	\question First question?
	
		\begin{solution}[1in]{}
		\end{solution}
	
	\question Second question?
	
		\begin{solution}{}
		
			Here's a solution environment that doesn't have a height set, but does have content.
			
			(The previous solution environment has no height and no content.)		
		
		\end{solution}
\end{ctqs}

		
\begin{infobox}
	Here is some useful information that might help the students with the next few critical thinking questions.
	In some cases, it might include an equation:
	\begin{equation*}
		a^2 + b^2 = c^2
	\end{equation*}
	or even more than one equation:
	\begin{equation*}
		x=\frac{-b \pm \sqrt{b^2-4ac}}{2a}
	\end{equation*}
\end{infobox}

\begin{ctqs}

	\question First question?
		\begin{solution}[1in]{}
		\end{solution}
	
	\question Second question?
		\begin{solution}[1in]{}
		\end{solution}
\end{ctqs}



\begin{exercises}

	\exercise First exercise ~\answer{a short answer}
	\exercise Second exercise
	
\end{exercises}


\begin{problems}

	\problem First exercise
	\problem Second exercise
	
\end{problems}


	
\end{activity}